\chapter{Documenting your TUI app}
\label{ch:docs}

\begin{aside}
A TUI without docs feels mysterious. The keys are the
documentation; the rest is footnotes.
\end{aside}

\section{The help screen}

The first stop. F1 or \code{?}. Every action visible.
See Chapter~\ref{ch:help}.

\section{The man page}

A traditional Unix tool ships a man page. Generated from
markdown or asciidoc.

\section{The README}

Quick install, quick run, a screenshot or asciinema cast.

\section{Tutorials}

Step-by-step walkthroughs for common workflows. Either in
the help screen (interactive) or as separate web docs.

\section{Reference}

Every key, every command, every option. Generated from
the keymap and config schema where possible.

\section{Changelog}

Per release: what's new, what changed, what broke.

\section{Examples}

Real config files, example workflows, screenshots of
common screens.

\section{Recap}

\begin{recap}
\begin{itemize}[leftmargin=*]
  \item Help screen + man + README cover most users.
  \item Reference for the determined.
  \item Examples for the visual learners.
  \item Changelog for upgrade fears.
\end{itemize}
\end{recap}

\section*{Workshop}

\begin{enumerate}[leftmargin=*]
  \item Write a one-paragraph README intro.
  \item Build an in-app help screen.
  \item Record an asciinema cast of your app.
\end{enumerate}
